\begin{afterword}
\summaryhead

The post opens with the Talmud's ``Whoever saves a single life, it is as if he had saved the whole world,'' and asks the reader to feel its warm glow. The author reports that turning one person's life around with an anonymous blog comment felt as good as a major scientific insight, and suggests that saving one life probably feels as good as saving the world.

Beyond the feeling, the post argues, the difference is gigantic. It names three views on which saving more lives need not be proportionally better, and holds against them that two lives are twice as valuable as one. A rescuer who pulls one child off a railroad track and lets the second die, being already ``unimaginably'' ahead, cannot be right, and a philanthropist choosing between cures is in the same position. When lives are at stake the duty is to maximize, and whoever knowingly saves fewer is as damned as a murderer. An addendum notes that obvious ways of spending money to save lives often fail.

\responsehead

The post's conclusion, that one should save the greater number, is widely shared, and not only by consequentialists. The Stanford Encyclopedia reports that many contractualists also want to capture the intuition that one should save five rather than one. The opposing view the post describes first has a published defender: John Taurek argued in 1977 that no bystander is required to save five rather than one, and proposed flipping a coin. So the post takes a side in a real dispute.

The argument, though, does not address the point in dispute. The railroad case works against the third view read literally, since if one life's value is already beyond measure, a second adds nothing. But the first two views also condemn the rescuer, because the second child can be saved at no cost to the first. The disagreement with a view like Taurek's concerns choices where saving some means losing others, as in the philanthropist's case. The post moves to that case with ``Why should it be any different'' and answers ``I don't think it is different,'' without further argument.

The last sentence goes further still. That the duty to maximize has ``the same strength'' as the duty to save lives, and that choosing fewer damns one ``as thoroughly as any murderer,'' equates letting die with killing. That is a long-disputed claim which, on the Stanford Encyclopedia's account, anti-consequentialists almost universally reject. The post asserts it.

The evidence about feeling is the author's own experience, and it covers one side of the comparison; the claim about saving the world is offered with ``probably.'' The epigraph's wording, ``a single life,'' is an attested reading of the Mishnah, though the printed editions add ``from Israel.''

In short: a widely shared conclusion, that one should save the greater number, argued from a case with no trade-off and applied to a case with one; its last step, that choosing fewer damns one ``as thoroughly as any murderer,'' is asserted without argument.
\end{afterword}
